\chapter{Building a wing from parts; MIMO responses}\label{ch:hybrid}

Chapter~\ref{ch:wing} treated the Goland wing as a single object, with
one characteristic function and one tip transfer function. This chapter
addresses two further needs. The first is \emph{assembly from parts}:
describe each substructure by its ports and connect the parts, the idea
of the transfer-matrix and dynamic-stiffness methods for beams
\citep{banerjee2016}. The second is \emph{several inputs at once}: the
response to simultaneous loads, the analogue of \code{lsim} for a
multi-input, multi-output system. Nothing is discretized. Each block is
the exact solution of the beam equations with exact Theodorsen strip
loads; the matrices are small because they describe ports, not truncated
modes. The code is \path{examples/continuous_wing/hybrid_quickstart.m}.

\section{A strip as a hybrid two-port}

A strip of constant properties maps its left state to its right state
exactly, with displacements $q=[w,w',\alpha]$ and internal loads
$p=[V,M,T]$:
\begin{equation}
 \begin{bmatrix}q_R\\p_R\end{bmatrix}
 =E(s)\begin{bmatrix}q_L\\p_L\end{bmatrix},\qquad
 E=e^{A(s)\ell}=\begin{bmatrix}a&b\\c&d\end{bmatrix}.
\end{equation}
A neighbour imposes a displacement at one end and a load at the other.
The strip is therefore rewritten with mixed inputs: $q_L$ and the load
$f_R=p_R$ are given, and $f_L=-p_L$ and $q_R$ are computed. Solving
$p_R=cq_L+dp_L$ for $p_L$ gives
\begin{equation}\label{eq:hybridports}
 \begin{bmatrix}f_L\\q_R\end{bmatrix}
 =H(s)\begin{bmatrix}q_L\\f_R\end{bmatrix},\qquad
 H=\begin{bmatrix}d^{-1}c&-d^{-1}\\a-bd^{-1}c&bd^{-1}\end{bmatrix}.
\end{equation}
These are the hybrid (h-) parameters of a two-port, familiar from circuit
theory, with three displacements and three loads per port. $H$ is
$6\times6$ because the beam has six state variables, and its entries are
transcendental functions of $s$ (Figure~\ref{fig:hybridports}).

\begin{figure}[htbp]\centering
\begin{tikzpicture}[>=Latex,thick,font=\small,x=1cm,y=1cm,
  inp/.style={->,blue!70!black}, outp/.style={->,red!75!black}]
  % Two elements
  \foreach \x/\n in {0/A,8/B} {
    \filldraw[fill=teal!12,draw=teal!60!black,rounded corners] (\x,-0.9) rectangle (\x+3,0.9);
    \node at (\x+1.5,0.3) {element \n};
    \node[font=\scriptsize] at (\x+1.5,-0.35) {$\begin{bmatrix}f_L\\q_R\end{bmatrix}=H_\n(s)\begin{bmatrix}q_L\\f_R\end{bmatrix}$};
  }
  % Left external port of A
  \draw[inp]  (-1.6,0.45)--(0,0.45) node[pos=0.35,above] {$q_L$};
  \draw[outp] (0,-0.45)--(-1.6,-0.45) node[pos=0.65,below] {$f_L$};
  % Interface between A and B
  \draw[outp] (3,0.45)--(4.2,0.45) node[pos=0.5,above] {$q_R^A$};
  \draw[inp]  (4.2,-0.45)--(3,-0.45) node[pos=0.5,below] {$f_R^A$};
  \draw[inp]  (6.8,0.45)--(8,0.45) node[pos=0.5,above] {$q_L^B$};
  \draw[outp] (8,-0.45)--(6.8,-0.45) node[pos=0.5,below] {$f_L^B$};
  \node[draw,dashed,rounded corners,fill=white,align=center,font=\scriptsize,inner sep=3pt]
      at (5.5,0) {$q_R^A=q_L^B$\\$f_R^A+f_L^B=0$};
  \node[font=\scriptsize,align=center] at (5.5,-1.65) {interface: solved and eliminated\\(not a product $H_AH_B$)};
  % Right external port of B
  \draw[inp]  (12.6,-0.45)--(11,-0.45) node[pos=0.35,below] {$f_R$};
  \draw[outp] (11,0.45)--(12.6,0.45) node[pos=0.65,above] {$q_R$};
  % Legend
  \draw[inp] (0,-2.0)--(0.7,-2.0); \node[anchor=west,font=\scriptsize] at (0.75,-2.0) {input (given)};
  \draw[outp] (0,-2.45)--(0.7,-2.45); \node[anchor=west,font=\scriptsize] at (0.75,-2.45) {output (computed)};
  \node[font=\scriptsize,align=left,anchor=west] at (8.4,-2.2)
      {$q=[w,\,w',\,\alpha]$: displacements\\$f=[V,\,M,\,T]$: loads};
\end{tikzpicture}

\caption{Two hybrid elements. Inputs (blue) are given and outputs (red)
computed. The interface conditions are solved and eliminated.}
\label{fig:hybridports}
\end{figure}

\section{Connection is not multiplication}

At a junction, the displacements are shared and the mutual loads balance:
$q_R^A=q_L^B=q_I$ and $f_R^A+f_L^B=0$. With $H^A,H^B$ partitioned into
$3\times3$ blocks,
\begin{equation}
 (I+H^A_{22}H^B_{11})\,q_I=H^A_{21}\,q_L-H^A_{22}H^B_{12}\,f_R,
\end{equation}
and eliminating $q_I$ yields the hybrid matrix $H^{AB}$ of the pair, with
the same port structure (\code{wing\_hybrid\_join}). The product
$H^AH^B$ would feed the loads of $A$ into the displacement inputs of $B$
and is meaningless. Two joined half-span strips reproduce one full-span
strip to $3\times10^{-16}$. With a clamped root ($q_L=0$) and tip loads
$u=f_R$, the tip displacements are $q_R=H^{AB}_{22}u$: $H^{AB}_{22}$ is the
$3\times3$ tip transfer matrix.

The implicit assembly of Chapter~\ref{ch:wing} keeps the interface states
as unknowns and gives $K(s)x=Bu$, $y=Cx$, that is $G=CK^{-1}B$. The
structured model \code{cdyn} represents it, and \code{ceval} evaluates
all nine channels with one factorization of $K$. The two assemblies
impose the same conditions and differ only in when the interface unknowns
are eliminated. At $s=3+70i$ and 120~m/s they agree to $2\times10^{-15}$.

\section{Artificial poles of the hybrid form}

The hybrid form inverts $d$ and $I+H^A_{22}H^B_{11}$, which depend on
$s$. $d$ is singular exactly when the piece alone, clamped at its left
end and free at its right end, has a mode. That is a property of the
piece, not of the assembled wing. In vacuum, a half-span piece has such a
mode at $4\times49.49=197.97$~rad/s. The whole wing has none there (its
nearest mode is at 261~rad/s), and its implicit transfer is finite, while
the hybrid matrix of the piece does not exist. These artificial poles are
the reason why the characteristic function must be built from the
implicit matrix: $\det H$ has poles that belong to no physical mode. The
functions report a singular pivot (\code{available=false}, or an error
from the transfer evaluators) and never regularize it. With air, the
artificial poles leave the imaginary axis; along $\operatorname{Re}s=6$ the
two assemblies agree to $1.6\times10^{-9}$ up to $2\times10^4$~rad/s.

\section{Matrix models and simultaneous inputs}

\code{cmimo} builds a transfer matrix, from channels, a constant matrix
or a matrix-valued handle; \code{cdyn} builds a structured model
$H(s)x=B(s)u$, $y=C(s)x+D(s)u$. \code{clsim}, \code{cstep},
\code{cimpulse} and \code{ckernel} accept both. \code{clsim} takes one
input column per input and returns one column per output, with all
inputs acting together. \code{cstep} and \code{cimpulse} return
$N_t\times n_y\times n_u$ arrays, one experiment per input. Each channel
uses the scalar inversion of Chapter~\ref{ch:response}. Evaluations of the
matrix model are shared between channels within one call, and a kernel
bank from \code{ckernel} stores the kernels of every channel.

Figure~\ref{fig:hybridsup} shows the tip response of the wing at
120~m/s to a force $1+0.2\sin10t$~kN and a torque $0.3\cos5t$~kN\,m. One
bank is prepared (about 30~s), and three \code{clsim} calls give the
response to both loads, to the force alone and to the torque alone. Each
call takes milliseconds and evaluates the wing model zero times. The
coupling is strong in both directions. The force twists the wing about as
much as the torque, and the torque raises the tip by up to 4~mm, partly
cancelling the deflection due to the force. The coupling is inertial
(centre of mass behind the elastic axis) and aerodynamic. The combined
response equals the sum of the two contributions to $9\times10^{-16}$. The inversion bound 5 is an assumption
supported by the root counts of Chapter~\ref{ch:wing}, not a proof.

\begin{figure}[htbp]\centering
\includegraphics[width=\textwidth]{hybrid_superposition.png}
\caption{Tip deflection and twist at 120~m/s: contributions of the force
and of the torque, and the response to both (black).}
\label{fig:hybridsup}
\end{figure}

Each output is a sum of channel contributions, and so is its error.
\code{clsim} adds the estimated channel errors of each output and
compares the total with $\mathrm{AbsTol}_i+\mathrm{RelTol}|y_i|$, with one
absolute tolerance per output in its own units. Destructive cancellation
between channels therefore shows up as an unresolved output instead of
being hidden.

\section{Verification}

\begin{itemize}
  \item \textbf{Transfer matrices.} Hybrid and implicit tip transfer
        matrices, all nine channels, agree to $1.6\times10^{-13}$. This
        holds for dry, uniform and tapered wings, with 1, 2, 4 and 8
        strips, at five complex frequencies.
  \item \textbf{Structured model.} \code{cdyn} agrees with the scalar
        tip transfer of Chapter~\ref{ch:wing} to $2\times10^{-14}$.
  \item \textbf{Time responses.} The hybrid and implicit kernel banks
        agree to $5\times10^{-14}$. The matrix bank agrees with the sum of
        independent scalar \code{clsim} calls within the requested output
        tolerance: $3\times10^{-4}$~mm against $10^{-2}$.
  \item \textbf{A rational oracle.} A $2\times2$ rational model with a
        different transport delay in each channel agrees with MATLAB's
        \code{lsim} with zero-order hold to $1.4\times10^{-9}$.
\end{itemize}
These checks show that the representations agree; the physical
references remain those of Chapter~\ref{ch:wing}.

\section{Scope}

The new functions provide matrix models and MIMO zero-state responses.
Pole and zero searches remain scalar: matrix characteristic values, MIMO
transfer poles and transmission zeros are not implemented, and the
characteristic function of a matrix problem is still an analytic
determinant such as $\det K$. The hybrid connector is an example for a
chain of beam strips, not a general network assembler.
